%-------------------------- model.tex -------------------------%
%Exemple per a usar l'estil tfm.class per a la preparaci\'o
%de monografies de treballs final de m\`aster.
%
%
%----------------------------- FMETFM ----------------------------%
\documentclass{fmetfm}

%%%%%%%%%%packages
\usepackage[catalan,spanish,english]{babel}
%%%%%%%%%%%%%%%%%%%%%%%%%%%%%%%%%%%%%%%%%
%permet l'�s d'un dels tres idiomes.
%es poden seleccionar amb l'opci\'o
%\selectlanguage que hi ha despr\'es
%de \begin{document}
%%%%%%%%%%%%%%%%%%%%%%%%%%%%%%%%%%%%%%%%
\usepackage[latin1]{inputenc}
%%%%%%%%%%%%%%%%%%%%%%%%%%%%%%%%%%%%%%%%
%permet l'�s d'accents, �,�, etc.
%%%%%%%%%%%%%%%%%%%%%%%%%%%%%%%%%%%%%%%%%
\usepackage{amsmath}%permet usar simbols matematics especials
\usepackage{amsfonts}%permet usar fonts matematiques especials
\usepackage{latexsym}%permet usar simbols matematics especials
\usepackage{epsfig}%permet incorporar grafics en .eps
\usepackage{psfig}%permet incoporar grafics en .ps
\usepackage{graphicx}%permet incorporar grafics
%\usepackage{makeidx}%per elaborar un index de termes
%%%%%%%%%%%%%%%%%%%%%%%%%%%%%%%%%%%%%%%%%
% Per als entorns matem\`atics feu servir:
% \begin{lemma}       ...  \end{lemma}
% \begin{proposition} ...  \end{proposition}
% \begin{theorem}     ...  \end{theorem}
% \begin{corollary}   ...  \end{corollary}
% \begin{conjecture}  ...  \end{conjecture}
% per a les demostracions:
% \begin{proof}       ...  \qed \end{proof}
% i per  a les notes o observacions:
% \begin{note}    ... \end{note}
% \begin{remark}  ... \end{remark}
%%%%%%%%%%%%%%%%%%%%%%%%%%%%%%%%%%%%%%%%%%%%
\makeindex

%

%--------------------------------------------------------------%
%                  Espai pels autors.                   %
%--------------------------------------------------------------%
%%-----------------------------------------------------------------
%% P�gina inicial
%%-----------------------------------------------------------------
%\title{Exemple d'�s de la classe `fmetfm' per a l'elaboraci�
%de mem�ries de treballs final de m�ster}
%
%\treball{Tipus de treball: Tesi de m�ster, Treball Fi de Carrera}
%\advisor{Nom del director del treball}
%\author{Nom de l'autor}
%
%\treball{Tipus de treball: Tesi de m�ster, Treball Fi de Carrera}
%\advisor{Nom del director del treball}
%
%
%
%%-------------------------------------------------------------%
\begin{document}

% Seleccioneu l'idioma:
\selectlanguage{catalan}
%\selectlanguage{spanish}
%\selectlanguage{english}
%----------------------------------------------------------------
% Inici del document
%------------------------------------------------------------------

%-----------------------------------------------------------------
% P�gina inicial
%-----------------------------------------------------------------
\title{T�tol del treball}


\def\autor{Nom de l'autor}

\def\treball{Tipus de treball\\ (Tesi de m�ster, Projecte fi de carerra,...)}

\def\advisor{Nom(s) del director(s)}

\vfill

\def\departament{Nom del Departament}



%-------------------------------------------------------------%
\maketitle



%-------------------------------------------------------------%
%%%Dedicatoria (opcional)
%-------------------------------------------------------------%
\dedicatoria{Com que la dedicat�ria normalment �s una frase breu que
vol afectar tot el contingut de l'obra, �s convenient que aparegui
just despr�s de la portada i en una p�gina a part. S'acostuma a
col.locar al marge superior dret del full}
%-------------------------------------------------------------%
%-------------------------------------------------------------%
%%%Prefaci (opcional)
%-------------------------------------------------------------%

\begin{preface}
\thispagestyle{empty}
 Tot i que no �s obligat�ria la seva pres�ncia,
el pr�leg i el prefaci s�n dos apartats que normalment apareixen a
les obres. La funci� de tots dos �s transmetre opinions o reflexions
sobre l'obra, explicar el proc�s d'elaboraci� o argumentar-ne
l'oportunitat. El pr�leg el redacta una persona diferent de l'autor,
i el prefaci, l'autor mateix.
\end{preface}


\clearpage

%---------------------------------------------------------------
% Abstract: un resum del contingut del treball (sobre 500 paraules)
% Ha de contenir una relaci� de paraules clau i dels codis MSC
%(Math. Subject Classification 2000)
% Ha d'estar tamb� en angl�s.
%---------------------------------------------------------------
\pagenumbering{roman}
\setcounter{page}{5}
\markboth{}{}

\begin{abstract}

\keywords{Paraules clau} \msc{Codis de la {\it Mathematic Subject
Classification}}

\vspace{1cm}
%%%Aqui va el text

Un resum del contingut del treball d'unes 500 paraules. El resum
d�na una idea global del treball i en descriu les aportacions m�s
significatives.

\end{abstract}

\newpage

%%%Si el text no es en angles afegir una versi� anglesa

\begin{abstracteng}

\keywordseng{Here the keywords}
\msc2000{Mathematical Subject
Classification}

\vspace{1cm}
%%%Here the text

An abstract of the contents of the work. About 500 words. It gives a
global description of the work and its main contributions.

\end{abstracteng}

\cleardoublepage

%%%%%%%%%%%%%%%%%%%%%%%%%%%%%%%%%%%%%%%%%%%%%%%%%%%%%%%%%%
%%Notation (optional)
%%%%%%%%%%%%%%%%%%%%%%%%%%%%%%%%%%%%%%%%%%%%%%%%%%%%%%%%%%
\begin{notation}
\begin{tabular}{cl}
${\mathbb N}$ & Nombres naturals\\[0.3em]
${\mathbb Z}$ & Nombres enters\\[0.3em]
${\mathbb Q}$ & Nombres racionals\\[0.3em]
${\mathbb R}$ & Nombres reals\\[0.3em]
${\mathbb C}$ & Nombres complexos
\end{tabular}
\end{notation}




\cleardoublepage
%--------------------------------------------------------------
% Taula de continguts
% S'elabora autom�ticament
%-------------------------------------------------------------
 \tableofcontents
 \cleardoublepage
%------------------------------------------------------------
% Cos del treball
%Es poden fer servir Parts, Capitols, Seccions i Subseccions
%-------------------------------------------------------------
\pagenumbering{arabic}

\input{Introduction}
\part{Descripci� de l'estructura}
\input{cap1}
\input{cap2}
\part{Elements especials}
\input{cap3}
\appendix{Fora del cos} Els ap�ndixs poden contenir demostracions
especialment t�cniques, prgrames d'ordinador, descripci�
d'algorismes, etc.


\printindex

\begin{thebibliography}{99}
\bibitem{pdlatex} D.P. Carlsihe, S.P.Q. Ratz, The graphicx package, {\it Latex documentation} (1999).
\bibitem{index} David M. Jones, A new implementation of LATEX's indexing
commands {\it Latex documentation} (1995)
\bibitem{latex} Leslie Lamport, LaTeX: A Document Preparation System (2nd Edition). Addison-Wesley Series on Tools and Techniques for Computer
T
\end{thebibliography}

%\cleardoublepage \sloppy
%\begin{raggedright}
%\input{exemple_tfm.ind}
%\end{raggedright}



%--------------------------------------------------------------
%Seccions addicionals: Ap�ndixs, Glosari, Bibliografia
%
%Els ap�ndixs son com capitols pero van precedits de \apendix,
%es a dir: \apendix \chapter{Nom de l'apendix}
%
%\input{apendix1}

%Per a la bibliografia es pot fer servir els fitxers .bib
%(consulteu un manual de LaTex per al seu us)
%o simplement amb \begin{thebibliography}{100} \end{thebibliography}
%
%\bibliographystyle{amsplain}
%\bibliography{tfm}

%El glosari de termes es crea automaticament. Consulteu un manual
%de LaTex

%--------------------------------------------------------------



%
\end{document}
%--------------------------------------------------------------%
%--------------------------------------------------------------%
